\chapter{Statistics and Regression}
\label{ch:regression}

Many causal estimators studied later can be represented as linear regressions
or as solutions to moment equations.  This chapter develops least squares,
its large-sample properties, heteroskedasticity-robust standard errors, tests,
and confidence intervals.

\section{Identification, estimation, and inference}

From a frequentist perspective, data are generated by a fixed but unknown
process.  For example,
\[
  Y=g(X,\varepsilon),
\]
together with the distribution of $(X,\varepsilon)$, describes a data
generating process (DGP).  We observe a sample of $(Y_i,X_i)$ but not the
structural disturbance $\varepsilon_i$ or the function $g$.

Suppose the target is a feature $\theta$ of the DGP.  Econometric analysis has
three distinct tasks:
\begin{description}
  \item[Identification.] Does the population distribution determine a unique
        value of $\theta$ under the maintained assumptions?
  \item[Estimation.] Can we construct a sample rule $\widehat\theta$ that is
        close to $\theta$?
  \item[Inference.] How much sampling uncertainty surrounds
        $\widehat\theta$?  Tests and confidence intervals answer related but
        different versions of this question.
\end{description}
No amount of data repairs a failure of identification: if two parameter
values produce the same observable distribution, the sample cannot tell them
apart.

\section{The linear model and least squares}

Consider an i.i.d. sample satisfying the population linear projection
\begin{equation}
\label{eq:linear-model}
  Y_i=X_i'\beta+\varepsilon_i,
  \qquad \mathbb E(X_i\varepsilon_i)=0,
\end{equation}
where $X_i$ is $k\times1$ and usually contains a constant.  Stacking
observations gives
\[
  Y=X\beta+\varepsilon,
\]
where $X$ is $n\times k$ and row $i$ is $X_i'$.  At this stage, $\beta$ need
not have a causal interpretation; later designs supply the assumptions that
make particular regression coefficients causal.

\subsection{A method-of-moments estimator}

The orthogonality condition in \eqref{eq:linear-model} implies
\[
  \mathbb E(X_iY_i)=\mathbb E(X_iX_i')\beta.
\]
If $Q=\mathbb E(X_iX_i')$ is nonsingular, then the population distribution
identifies
\begin{equation}
\label{eq:beta-population}
  \beta=Q^{-1}\mathbb E(X_iY_i).
\end{equation}
Nonsingularity rules out perfect multicollinearity.  For example, a constant,
a male indicator, and a female indicator cannot all be included when the two
indicators always sum to one.  One of the three columns must be removed.

The sample analogue replaces expectations by sample averages:
\begin{align}
\widehat\beta
  &=\left(\frac1n\sum_{i=1}^{n}X_iX_i'\right)^{-1}
    \left(\frac1n\sum_{i=1}^{n}X_iY_i\right) \notag\\
  &=(X'X)^{-1}X'Y. \label{eq:ols}
\end{align}
This is both a method-of-moments estimator and the ordinary least-squares
(OLS) estimator.  The least-squares interpretation follows because
\eqref{eq:ols} minimizes
$\sum_i(Y_i-X_i'b)^2$ when $X$ has full column rank.

Method of moments is a broad and useful recipe:
\begin{enumerate}
  \item obtain population moment restrictions;
  \item replace population moments by sample averages;
  \item solve the sample moment equations.
\end{enumerate}
It does not eliminate the hard work: credible moment restrictions still come
from economic reasoning and an identification argument.

\begin{rcode}
set.seed(5280)
n <- 100
x <- rnorm(n)
e <- rnorm(n)
y <- 0.5 + 2 * x + e
X <- cbind(1, x)

# Matrix formula
bhat <- solve(crossprod(X), crossprod(X, y))
bhat

# R's regression routine
coef(lm(y ~ x))
\end{rcode}

\subsection{Fitted values and residuals}

The fitted values and residuals are
\begin{align*}
  \widehat Y&=X\widehat\beta
    =X(X'X)^{-1}X'Y\equiv PY,\\
  \widehat\varepsilon&=Y-\widehat Y
    =\{I-X(X'X)^{-1}X'\}Y\equiv MY.
\end{align*}
The matrices $P$ and $M=I-P$ project $Y$ respectively onto the column space
of $X$ and its orthogonal complement.  The normal equations imply
$X'\widehat\varepsilon=0$.

\section{Statistical properties of OLS}

\subsection{Unbiasedness}

\begin{definition}[Unbiasedness]
An estimator $\widehat\beta$ is unbiased for $\beta$ if
$\mathbb E(\widehat\beta)=\beta$.
\end{definition}

Unbiasedness is a finite-sample property.  For conditional unbiasedness of
OLS, assume
\begin{equation}
\label{eq:zero-conditional-mean}
  \mathbb E(\varepsilon\mid X)=0.
\end{equation}
Under random sampling, the familiar condition
$\mathbb E(\varepsilon_i\mid X_i)=0$ implies
\eqref{eq:zero-conditional-mean}.  Conditional on a full-rank design matrix,
\begin{align*}
  \mathbb E(\widehat\beta\mid X)
  &=\mathbb E\{\beta+(X'X)^{-1}X'\varepsilon\mid X\}\\
  &=\beta+(X'X)^{-1}X'\mathbb E(\varepsilon\mid X)\\
  &=\beta.
\end{align*}
The law of iterated expectations therefore gives
$\mathbb E(\widehat\beta)=\beta$.  Notice that the target is $\beta$, not
zero.

\begin{rcode}
set.seed(5280)
n <- 100
B <- 500
bhat <- replicate(B, {
  x <- rnorm(n)
  y <- 0.5 + 2 * x + rnorm(n)
  coef(lm(y ~ x))
})
rowMeans(bhat)
\end{rcode}

\subsection{Consistency}

Using $Y=X\beta+\varepsilon$,
\begin{equation}
\label{eq:ols-error}
  \widehat\beta-\beta
  =\left(\frac1n\sum_iX_iX_i'\right)^{-1}
   \left(\frac1n\sum_iX_i\varepsilon_i\right).
\end{equation}
Suppose the relevant moments exist, observations are i.i.d.,
$Q=\mathbb E(X_iX_i')$ is nonsingular, and
$\mathbb E(X_i\varepsilon_i)=0$.  The WLLN and continuous mapping theorem
give
\[
  \frac1n\sum_iX_iX_i'\pto Q,
  \qquad
  \frac1n\sum_iX_i\varepsilon_i\pto0,
\]
and hence $\widehat\beta\pto\beta$.  Conditional mean zero is sufficient but
not necessary for this conclusion; the weaker orthogonality condition is
enough.

\subsection{Asymptotic normality}

Multiply \eqref{eq:ols-error} by $\sqrt n$:
\[
  \sqrt n(\widehat\beta-\beta)
  =\left(\frac1n\sum_iX_iX_i'\right)^{-1}
   \left(\frac1{\sqrt n}\sum_iX_i\varepsilon_i\right).
\]
Let
\[
  Q=\mathbb E(X_iX_i'),
  \qquad
  \Omega=\mathbb E(\varepsilon_i^2X_iX_i').
\]
Under the preceding conditions and finite second moments for
$X_i\varepsilon_i$, the multivariate CLT and Slutsky's theorem yield
\begin{equation}
\label{eq:ols-an}
  \sqrt n(\widehat\beta-\beta)
  \dto N(0,\Sigma),
  \qquad
  \Sigma=Q^{-1}\Omega Q^{-1}.
\end{equation}
The matrix $Q^{-1}\Omega Q^{-1}$ is called a sandwich variance: $Q^{-1}$ is
the bread and $\Omega$ is the filling.  Equivalently, for large $n$,
$\widehat\beta$ is approximately $N(\beta,\Sigma/n)$.

\subsection{Heteroskedasticity-robust standard errors}

Neither consistency nor \eqref{eq:ols-an} requires constant conditional error
variance.  A heteroskedasticity-robust estimator of $\Sigma$ is
\begin{equation}
\label{eq:robust-av}
  \widehat\Sigma
  =\widehat Q^{-1}\widehat\Omega\widehat Q^{-1},
  \quad
  \widehat Q=\frac1n\sum_iX_iX_i',
  \quad
  \widehat\Omega=\frac1n\sum_i
       \widehat\varepsilon_i^2X_iX_i'.
\end{equation}
If $\widehat\Sigma_{jj}$ is diagonal element $j$, the standard error of
$\widehat\beta_j$ is
\[
  \operatorname{se}(\widehat\beta_j)
  =\sqrt{\widehat\Sigma_{jj}/n}.
\]
It estimates the standard deviation of $\widehat\beta_j$, not the standard
deviation of the limiting variable in \eqref{eq:ols-an}.

\section{Testing}

Consider
\[
  H_0:\beta_j=\beta_j^0
  \qquad\text{against}\qquad
  H_1:\beta_j\neq\beta_j^0.
\]
A Type I error rejects a true null; a Type II error fails to reject a false
null.  A level-$\alpha$ procedure controls the probability of a Type I error
at approximately $\alpha$.  Its power is the probability of rejection when
the alternative is true.

Under $H_0$, the studentized statistic
\begin{equation}
\label{eq:t-stat}
  T_n=\frac{\widehat\beta_j-\beta_j^0}
  {\operatorname{se}(\widehat\beta_j)}
  \dto N(0,1).
\end{equation}
For a two-sided asymptotic level-$\alpha$ test, reject when
$|T_n|>z_{1-\alpha/2}$.  Common critical values are 1.645 for 10 percent,
1.960 for 5 percent, and 2.576 for 1 percent tests.

The two-sided $p$-value is
\[
  2\{1-\Phi(|t_{\mathrm{obs}}|)\}.
\]
It is the null probability of observing a statistic at least as extreme as
the realized one.  It is \emph{not} the probability that the null hypothesis
is true.  Failure to reject should not be described as acceptance: it may
reflect low power.

Because standard errors shrink with sample size, a very small effect can be
statistically significant in a very large sample.  Statistical significance
must therefore be distinguished from economic importance.

\section{Confidence intervals}

An asymptotic $100(1-\alpha)$ percent confidence interval is
\begin{equation}
\label{eq:ci}
  \left[
  \widehat\beta_j-z_{1-\alpha/2}\operatorname{se}(\widehat\beta_j),
  \widehat\beta_j+z_{1-\alpha/2}\operatorname{se}(\widehat\beta_j)
  \right].
\end{equation}
Its coverage probability approaches $1-\alpha$ under the maintained
assumptions.  The parameter is fixed; across repeated samples, the random
interval changes.  A 95 percent procedure covers the fixed true parameter in
approximately 95 percent of repeated samples.  It does not assign a 95
percent probability to the parameter after a particular frequentist interval
has been observed.

The interval and the two-sided test are dual: a candidate value is rejected
at level $\alpha$ exactly when it lies outside the corresponding confidence
interval (up to numerical and approximation differences).  Reporting point
estimates with standard errors or confidence intervals conveys both magnitude
and uncertainty.

\begin{rcode}
library(sandwich)
library(lmtest)

set.seed(5280)
n <- 500
x1 <- rnorm(n)
x2 <- rnorm(n)
y <- 0.5 + x1 + x2 + rnorm(n, sd = 0.5 + abs(x1))
fit <- lm(y ~ x1 + x2)

# HC1 heteroskedasticity-robust covariance matrix and coefficient table
V_HC1 <- vcovHC(fit, type = "HC1")
coeftest(fit, vcov. = V_HC1)

# A 95% robust confidence interval for the coefficient on x1
b1 <- coef(fit)["x1"]
se1 <- sqrt(V_HC1["x1", "x1"])
b1 + c(-1, 1) * qnorm(0.975) * se1
\end{rcode}

\section{Regression in causal designs}

We will almost never assume without argument that a regression coefficient is
a causal effect.  Instead, experiments or observational research designs will
produce identifying comparisons, and regression will often provide a
convenient way to compute them.  Always separate the statistical question
``How is this coefficient estimated?'' from the causal question ``Why does
this coefficient equal the effect of interest?''
